\chapter{Typical procedures}
\label{ch:workflow_tipici}
\index{Workflow!typical}\index{Procedures}\index{Use cases}

\epigraph{``To learn to do a new thing well, it helps to repeat
it five times. The first time gropingly, the second with caution,
the third with ease, the fourth with grace, the fifth with
mastery.''}{Adapted from \textsc{Beatrice Bizzarri},
\emph{Teaching notebook}, 1923}

\lettrine[lines=3]{T}{his} chapter gathers the six recurring
procedures from the experience of a timetable coordinator. Each
one is described as a sequence of numbered steps, with the tab to
open, the buttons to press, the expected effect. They are the
operational cookbook: to learn a school of thirty-two classes by
heart in two hours. The procedures are independent of one another:
a school that opens the timetable in August will follow the first
two; a school that steps in mid-year with an Excel file will jump
to the third.

\grokfig{fig_happy_path_aula}{The classroom at the first working
timetable: fifteen happy students and the weekly calendar on the
teacher's desk.}

\section{Procedure 1: a school from scratch}
\label{sec:wf_da_zero}

Starting point: empty database, a school to be configured for the
first time. Destination: a complete weekly timetable displayed.

\subsection{Step 1 -- configure the time grid}

Go to the \emph{Ore} (Hours) tab. Check that the working days and
the time slots reflect the reality of the school. The default
(Mon--Sat, 8:00--14:00, six slots) is fine for the vast majority
of Italian schools. Change it only if necessary. Press
\emph{Salva} (Save).

\paragraph{Attention.} Fix the grid now, not later. Changing it
afterwards invalidates the unavailability matrices already
entered.

\subsection{Step 2 -- load the registry data}

You have two roads.

\paragraph{Road A: import from Excel.}
Go to the \emph{Import bulk} tab. Download the templates of the
seven sheets (\texttt{teachers}, \texttt{subjects},
\texttt{classes}, \texttt{classrooms}, \texttt{curricula},
\texttt{students}, \texttt{groups}). Fill in each sheet. Re-import
in the same order, in \texttt{upsert} mode.

\paragraph{Road B: manual entry.}
Visit the tabs \emph{Indirizzi} (Tracks), \emph{Materie}
(Subjects), \emph{Docenti} (Teachers), \emph{Classi} (Classes),
\emph{Aule} (Classrooms), \emph{Studenti} (Students),
\emph{Gruppi} (Groups) in this order, and use the \emph{+ Nuovo}
(+ New) button of each to insert the rows. Longer, but suited to
small schools or to when a starting Excel sheet is missing.

\subsection{Step 3 -- assign the cattedre}

Go to the \emph{Cattedre} (Teaching posts) tab. For each class
you see a card with all the subjects of its track. Press
\emph{cambia} (change) next to each subject and choose the
teacher. For the first session it is best to start with the most
constrained teachers (example: those with a single possible
class), and finish with the \emph{jolly} (flexible, wildcard)
ones. Once configuration is complete, check with the
\emph{Warnings cattedre} button that no teacher is
\texttt{SFORA} (over quota) or \texttt{SOTTO} (under quota).

\subsection{Step 4 -- add unavailabilities and constraints}

For each teacher open the card from the \emph{Docenti} (Teachers)
tab and colour their matrix. Use red (\textsc{Hard}) for the real
unavailabilities; yellow (\textsc{Soft}) for the negative
preferences; blue (\textsc{Preferred}) for positively preferred
days.

For the logical constraints (e.g.\ ``Bianchi does not work on
Saturday'', ``in 5A no history on Friday'') go to the
\emph{Vincoli} (Constraints) tab, press \emph{+ Nuovo vincolo DSL}
(+ New DSL constraint), write the expression, \emph{Valida}
(Validate), \emph{Salva} (Save).

\begin{esempiobox}[title={Recurring constraints, in DSL}]
\begin{lstlisting}
# Bianchi does not work on Saturday
forall l in lessons where l.teacher == Bianchi:
    l.day != sab

# In 5A no history on Friday
forall l in lessons where l.class == 5A and l.subject == Storia:
    l.day != ven

# The gym is occupied on Tuesday morning
forall l in lessons where l.classroom == Palestra:
    l.day != mar or l.hour > 11
\end{lstlisting}
\end{esempiobox}

\subsection{Step 5 -- snapshot before the computation}

Go to the \emph{Dashboard}, \emph{Import / Export DB} card, press
\emph{Salva snapshot} (Save snapshot). This way you have a
guaranteed way back if the first pipeline produces strange
results.

\subsection{Step 6 -- pre-check}

On the \emph{Diagnostica} (Diagnostics) tab press
\emph{Hall pre-check}. The system verifies that the required hours
are formally serviceable by the available teachers. If the check
fails, the panel indicates the subset of classes and subjects in
deficit: do not launch the pipeline, fix things first.

\subsection{Step 7 -- launch the pipeline}

Go to the \emph{Workflow} tab, press \emph{Pipeline completa}
(Full pipeline). The lower log panel shows the phases one after
the other (Phase A, Phase B per day, metaheuristics, classrooms).
Depending on the size of the school:

\begin{itemize}
\item \texttt{small} (8 classes): 30--90 seconds;
\item \texttt{medium} (15--20 classes): 2--5 minutes;
\item \texttt{huge} (40+ classes): 8--20 minutes.
\end{itemize}

When it finishes, the green banner \emph{Pipeline completed}
takes you directly to the \emph{Orario} (Timetable) tab.

\subsection{Step 8 -- view and refine}

On the \emph{Orario} (Timetable) tab, navigate the four views (by
class, teacher, classroom, slot). If you find a lesson out of
place, drag it to another slot: the system checks the Hard
constraints in real time. If all is well, export with the
\emph{Esporta XLSX} (Export XLSX) button.

\grokfig{fig_no_solution_ape}{The confused bee: what to do when no
solution exists and a constraint has to be relaxed.}

\section{Procedure 2: importing data from another school's Excel}
\label{sec:wf_import_excel}

A school that receives the registry Excel sheets from another
institute, or from a previous year.

\paragraph{Premise.} Not all sheets arrive in the format of our
template. A small renaming job before the import is advisable.

\begin{enumerate}
\item Open the official template from the \emph{Import bulk} tab
  (\emph{Scarica template} / Download template button for
  \texttt{teachers}). Open the \emph{Istruzioni} (Instructions)
  sheet: there the canonical names and the accepted synonyms are
  listed (e.g.\ \texttt{nome}, \texttt{teacher\_name},
  \texttt{cognome\_nome} are all accepted for the \texttt{name}
  field).
\item In your starting Excel, rename the column headers with the
  canonical names (or any one of the aliases). Extra columns are
  ignored.
\item Save the file. On the \emph{Import bulk} tab choose
  \texttt{teachers}, \texttt{upsert} mode, drag the file. Press
  \emph{Importa} (Import).
\item The report panel below the button lists the imported rows
  and the discarded ones. The discarded ones report the index
  and the reason for the error (missing field, out-of-range
  value, and so on).
\item Repeat for the other six sheets, in the same order.
\end{enumerate}

\begin{erroricomunibox}
\textbf{``The file has broken accents.''} The importer accepts
Excel xlsx (recommended) or CSV. If the CSV comes from an
Italian Excel, save it in \emph{CSV UTF-8 (delimited)} format;
otherwise the letter \`e becomes a sequence of symbols and the
records are discarded.

\textbf{``The classes are imported but have no track.''}
Probably the \texttt{curriculum} column was missing from the
file, or the value did not match any track in the database. Fix
the error by importing the tracks first and then the classes.
\end{erroricomunibox}

\grokfig{fig_merge_rivoletti}{Three little streams of water flow
together into the river \texttt{main} downstream of the merge.}

\section{Procedure 3: re-running the pipeline after a change of constraints}
\label{sec:wf_relancia}

Typical scenario: the teachers' assembly asks for a substantial
change (e.g.\ ``let us add a co-teaching in all the first-year
classes''). The pipeline has to be re-run.

\begin{enumerate}
\item Snapshot of the current solution, from the \emph{Dashboard},
  \emph{Import / Export DB} card, \emph{Salva snapshot} (Save
  snapshot) button. The snapshot serves to compare before/after.
\item Add the new constraints (matrix, DSL, co-teachings\dots).
\item On the \emph{Diagnostica} (Diagnostics) tab press
  \emph{Hall pre-check} to verify that the request is still
  formally feasible.
\item Go to the \emph{Workflow} tab, press \emph{Pipeline
  completa} (Full pipeline). If the first pipeline had taken N
  minutes, the second will take a similar time.
\item Compare. On the \emph{Runs} tab you will see the new run at
  the top. Click it for the detail, compare the SOFT score with
  the previous one.
\item If the new one is better: it is done. If worse: go to the
  previous snapshot and press \emph{Ripristina} (Restore).
\end{enumerate}

\section{Procedure 4: editing a single lesson by hand}
\label{sec:wf_edit_manuale}

Scenario: the timetable has been in production for a week, the
principal asks to move 2B's English from Friday to Thursday.

\begin{enumerate}
\item On the \emph{Orario} (Timetable) tab, choose the \emph{Per
  classe} (By class) view, select 2B.
\item Drag the English cell from Friday to Thursday. As you drag,
  the candidate Thursday cells become coloured: green if
  acceptable, yellow if they worsen the SOFT, red if they violate
  a Hard.
\item Drop it over a green slot. The timetable updates; the
  \emph{Cost} panel at the top changes.
\item If the classroom was not compatible with the new slot, the
  \emph{Aula da scegliere} (Classroom to choose) modal opens:
  select one of the green alternatives.
\item Export the new timetable (XLSX) to distribute it.
\end{enumerate}

\paragraph{Suggestion.} The manual change does not automatically
update the saved constraints. If the change is to become permanent
(\emph{2B's English is always on Thursday}), it is best to add the
corresponding constraint in the \emph{Vincoli} (Constraints) tab
so that the next pipeline respects it without having to repeat the
drag-and-drop.

\grokfig{fig_supplenza_volante}{The substitute teacher flies
towards the classroom with her folder under her arm.}

\section{Procedure 5: handling a daily substitution}
\label{sec:wf_supplenza}

Daily scenario: three teachers are absent, the substitutions for
today have to be placed.

\begin{enumerate}
\item Open the \emph{Assenze e supplenze} (Absences and
  substitutions) tab.
\item Click the header of the current day (example:
  \emph{Marted\`{\i}} / Tuesday). A modal opens with the list of
  teachers; check the absent ones, confirm.
\item The timetable cells of those teachers turn red: uncovered
  classes.
\item For each red cell, click on the cell. A two-column modal
  opens: \emph{uncovered classes} on the left,
  \emph{available teachers} on the right. The available teachers
  are ordered by \texttt{graduatoria\_score} (ranking score): the
  most suitable appear at the top.
\item Drag a teacher onto the uncovered class; the cell turns
  green.
\item If all the cells are green, the day is covered. Export the
  absence bulletin with the button at the top right; it will be
  printed and posted.
\end{enumerate}

\begin{esempiobox}[title=Real case]
Liceo Galileo, Tuesday, 7:45. Absent are Bianchi (mathematics,
4 hours), Rossi (history, 2 hours) and Verdi (English, 3 hours):
nine uncovered cells. The system proposes at the top of the lists
three substitutes consistent with each subject, ordered by score.
The coverage closes in twelve minutes.
\end{esempiobox}

\grokfig{fig_backup_snapshot}{The backup hourglass with dated
snapshots in place of the sand.}

\section{Procedure 6: exporting and archiving the timetable}
\label{sec:wf_export}

At the end of the year, or when a consolidated version is closed,
it is best to archive three things:

\begin{enumerate}
\item The \emph{timetable in XLSX}. Go to the \emph{Orario}
  (Timetable) tab, \emph{Per classe} (By class) view, press
  \emph{Esporta XLSX} (Export XLSX). You will get a file with one
  sheet per class, formatted like the classic form of Italian
  schools. Repeat for the \emph{Per docente} (By teacher) view if
  needed.
\item The \emph{complete database}. Go to the \emph{Dashboard},
  \emph{Import / Export DB} card, press \emph{Esporta DB completo}
  (Export full DB). You will get a zip containing the raw SQLite
  file, the per-table CSVs and the metadata. It is the official
  backup.
\item The \emph{constraints}. On the same \emph{Import / Export
  constraints} card, press \emph{Esporta JSON} (Export JSON) (or
  xlsx). You will have a file with all the logical constraints
  written by the user, handy for reuse in similar schools.
\end{enumerate}

\paragraph{Archiving note.} It is best to save the three files in
a folder named \texttt{<school\_name>\_<school\_year>} on the
school's shared drive. Year after year, the archive becomes the
basis for historical studies (e.g.\ ``where did we get
worse/better'') and for restarting the next year from a known
structure.

\fleuron

\begin{biblioparvabox}
The six procedures in this chapter cover ninety per cent of the
real use cases. The residual situations (advanced diagnostics,
deep infeasibility, branch-and-price for enormous schools) are
treated in the specialist chapters:
\cref{ch:diagnostica_statistica_en} for statistical diagnostics,
\cref{ch:tecniche_avanzate_en} for the advanced techniques,
\cref{ch:lessons_learned_en} for the problematic scenarios
encountered in the pilot schools.
\end{biblioparvabox}
